本研究提出 SEER (Sub-question Entropy-driven Efficient Retrieval)，一個基於子問題層級的自適應檢索增強生成 (Retrieval-Augmented Generation, RAG) 框架。現有 RAG 系統在面對不同複雜度的查詢時，往往採用一體適用的檢索策略：對簡單問題進行不必要的檢索會引入雜訊並增加成本；對複雜的多跳推理問題僅進行單次檢索則經常導致答案不完整或錯誤。

近期的自適應 RAG 方法雖試圖解決此問題，但其決策粒度往往不適切：問題層級的路由過於粗糙，無法在推理過程中即時調整；詞元 (token) 層級的偵測則過於嘈雜。雖有方法採用子問題層級的決策 (如 DeepRAG)，但仰賴二元樹搜尋、SFT 與 DPO 等繁複訓練流程，犧牲了可攜性與可解釋性。

本研究的核心想法是：將複雜問題分解為原子化子問題後，針對每個子問題答案中的「首實體跨度 (first-entity span)」計算熵值，作為模型自信度的指標，再以此決定是否觸發檢索。SEER 包含三階段：(1) 資料建構，整合 MuSiQue、HotpotQA 與 2WikiMultiHopQA 資料集並進行品質過濾；(2) 以 SFT 訓練輕量化分解器；(3) 訓練無關的熵閾值門控機制。我們預期 SEER 在多跳問答基準上能達到與 DeepRAG 相當的 Exact Match 表現，同時顯著降低訓練成本並提升可解釋性。

\bigskip
\noindent\textbf{關鍵字：}檢索增強生成、自適應檢索、多跳問答、熵、子問題分解、大型語言模型
